\documentclass[border=10pt]{standalone}
\usepackage{fontspec}
\setmainfont{Noto Serif}[Ligatures=TeXOff]
\XeTeXgenerateactualtext=1
\newfontfamily\bengalifont{Noto Serif Bengali}[Script=Bengali, NFSSFamily=bnfam]
\usepackage[Latin,Bengali]{ucharclasses}
% group-free font switching (as in header-common.tex): robust inside braces, cells and quotes
\makeatletter
\def\BN@saved{\rmdefault}
\def\BN@bn{bnfam}
\newcommand\BNin{\edef\BN@tmp{\f@family}\ifx\BN@tmp\BN@bn\else\let\BN@saved\BN@tmp\fi\fontfamily{bnfam}\selectfont}
\newcommand\BNout{\fontfamily{\BN@saved}\selectfont}
\makeatother
\setTransitionsFor{Bengali}{\BNin}{\BNout}
% Typographic quotes and dashes are not in a ucharclasses class (class 0), so a Bangla word right
% after “ or ‘ would not trigger the Bangla font. Treat them as Latin so that every switch fires.
\XeTeXcharclass"2018=\XeTeXcharclass`A
\XeTeXcharclass"2019=\XeTeXcharclass`A
\XeTeXcharclass"201C=\XeTeXcharclass`A
\XeTeXcharclass"201D=\XeTeXcharclass`A
\XeTeXcharclass"2013=\XeTeXcharclass`A
\XeTeXcharclass"2014=\XeTeXcharclass`A
\XeTeXcharclass"2026=\XeTeXcharclass`A
\usepackage{newunicodechar}
\newunicodechar{।}{{\bengalifont\char"0964}}
\newunicodechar{॥}{{\bengalifont\char"0965}}

\usepackage[table]{xcolor}
\usepackage{booktabs,tabularx,array}
\definecolor{bapilrule}{HTML}{1F2937}
\definecolor{bapilstripe}{HTML}{F3F4F6}
\renewcommand{\arraystretch}{1.28}
\setlength{\tabcolsep}{5pt}

\begin{document}
\begin{minipage}{22cm}
{\large\bfseries সারণি ৭১।}\ {\large সমন্বিত প্রাতিষ্ঠানিক ও আর্থিক পথনকশা, ২০২৭ থেকে ২০৪১}\par\vspace{7pt}
\small
\rowcolors{2}{bapilstripe}{white}
\begin{tabularx}{\linewidth}{@{}>{\hsize=1.274\hsize\raggedright\arraybackslash}X>{\hsize=1.221\hsize\raggedright\arraybackslash}X>{\hsize=1.381\hsize\raggedright\arraybackslash}X>{\hsize=0.478\hsize\raggedright\arraybackslash}X>{\hsize=0.531\hsize\raggedright\arraybackslash}X>{\hsize=1.115\hsize\raggedright\arraybackslash}X@{}}
\arrayrulecolor{bapilrule}\toprule
\rowcolor{white}\textbf{পর্যায় ও সময়কাল} & \textbf{প্রাতিষ্ঠানিক রূপ ও নিয়োগকর্তা} & \textbf{সরকারি সংস্থা ও তাদের ভূমিকা} & \textbf{প্রত্যক্ষ পদ (পর্যায়ের শেষে)} & \textbf{বাজেট (কোটি টাকা)} & \textbf{পরবর্তী পর্যায়ের গেট} \\
\midrule
\textbf{ধাপ ০}, বছর ১-এর প্রায় ছয় মাস আগে (২০২৭-এর প্রথমার্ধ) & এখনো কোনো প্রতিষ্ঠান নেই; আইসিটি বিভাগ বা কোনো উন্নয়ন সহযোগীর নির্দেশে সমীক্ষা ও স্বাগতিকের আহ্বান & আইসিটি বিভাগ বা কোনো উন্নয়ন সহযোগী ধাপ ০-এর অর্থ দেয় ও ফোকাল পয়েন্ট মনোনীত করে (D5); রক্ষকেরা পুনর্ব্যবহারের শর্ত নিশ্চিত করে (D3); স্টিয়ারিং বোর্ডের মনোনয়ন প্রস্তুত হয় (D4) & নেই (চুক্তিভিত্তিক সমীক্ষা) & ০.৪৪ & স্বাগতিক নির্বাচিত; আইনি পর্যালোচনা ও ব্যয়-নিরূপণ সমীক্ষা সম্পন্ন; সিদ্ধান্ত D2 গৃহীত \\
\textbf{ধাপ ১-এর পাইলট}, বছর ১ (২০২৭-২৮ অর্থবছর) & সমঝোতা স্মারকের অধীনে বিশ্ববিদ্যালয়-ভিত্তিক ল্যাব (বিকল্প গ); স্বাগতিক প্রতিষ্ঠান নির্দিষ্ট মেয়াদের চুক্তিতে কর্মী নিয়োগ দেয় & আইসিটি বিভাগ সমঝোতা স্মারক স্বাক্ষর করে, স্টিয়ারিং বোর্ডের সহ-সভাপতিত্ব করে ও একজন ফোকাল পয়েন্ট মনোনীত করে; পরিকল্পনা কমিশন ও একটি ডেটা-রক্ষক সংস্থা বোর্ডে থাকে; রক্ষকেরা উন্মুক্ত নথিতে একসঙ্গে প্রবেশাধিকার অনুমোদন করে (D2 থেকে D4) & ৯ & ৪.৭৩ & মাস ১২-এর পর্যালোচনা এবং হালনাগাদ দ্বিতীয় ধাপের কর্মপরিকল্পনা ও বাজেট (D1.3) ধাপ ২ অনুমোদন করে ও কিস্তি ২ ছাড় করে \\
\textbf{ধাপ ২}, বছর ২ ও ৩ (২০২৮-২৯ থেকে ২০২৯-৩০ অর্থবছর) & একই রূপ; ERRC, TAB ও স্বচ্ছতা নিবন্ধন চালু & পাইলট সংস্থাগুলো ব্যবহারকারী ও ডেটা দেয়; স্বাগতিক সংস্থার সঙ্গে P1 ও P2-এর গ্রহণ-সিদ্ধান্ত (G4, মাস ২৪ থেকে); মাস ১৮-এ স্বাধীন মধ্যমেয়াদি মূল্যায়ন & ১৯ & ১৫.৭০ & ইতিবাচক চূড়ান্ত মূল্যায়ন (জুন ২০৩০); P1 ও P2 গেট G4 উত্তীর্ণ; কৌশলগত পর্যালোচনায় সম্প্রসারণের সিদ্ধান্ত (সারণি ৪৮) \\
\textbf{শর্তসাপেক্ষ জাতীয় সম্প্রসারণ}, বছর ৪ থেকে ৮ (২০৩০-৩১ থেকে ২০৩৪-৩৫ অর্থবছর) & বহুবর্ষী মূল অর্থায়নসহ ল্যাব; চারটি বিভাগীয় বিশ্ববিদ্যালয়-কেন্দ্র; কর্মীরা তখনো স্বাগতিক ও অংশীদার বিশ্ববিদ্যালয়ে নিযুক্ত & সরকারের মূল অবদান বাজেটের প্রায় অর্ধেকে পৌঁছায়; গৃহীত সরঞ্জাম সংস্থাগুলো শ্রেণি O সেবা হিসেবে পরিচালনা করে, নিযুক্ত বিশ্লেষকদের ধারণ করে ও সেবা-চুক্তি করে; ২০৩০-এর পরের জাতীয় এআই নীতিচক্রের সঙ্গে সামঞ্জস্য & ১১৮ & ২১৮.৮৫ & স্বাধীনতা, টেকসইতা, নীতিতে প্রয়োগ ও জন-আস্থার মানদণ্ডে ২০৩৫-এর কৌশলগত পর্যালোচনায় প্রাতিষ্ঠানিক রূপ নির্ধারণ \\
\textbf{প্রাতিষ্ঠানিকীকরণ}, ২০৩৫ থেকে & স্বাধীন জনস্বার্থমূলক ফাউন্ডেশন (বিকল্প ঘ), সরকারের সঙ্গে যুক্ত সংবিধিবদ্ধ কেন্দ্র, অথবা সংস্থাগুলো সেবা পরিচালনা করলে একটি ছোট গবেষণা-ল্যাব (অনুচ্ছেদ ৫.৯) & সংবিধিবদ্ধ কেন্দ্রের জন্য আইন বা সরকারের আনুষ্ঠানিক সিদ্ধান্ত প্রয়োজন (সংযুক্তি ন); তখন উন্মুক্ত নিয়োগের মাধ্যমে নতুন প্রতিষ্ঠানের চাকরিবিধির অধীনে পদগুলো নিয়মিত পদে পরিণত হতে পারে & ২০৩৫-এর পর্যালোচনায় নির্ধারিত & ২০৩৫-এর পর্যালোচনায় নির্ধারিত (ব্লুপ্রিন্ট ২.০ ও পরবর্তী) & ২০৪১-এ স্বাধীন পর্যালোচনা; জনমূল্য নিশ্চিত হলেই কেবল সনদ নবায়ন \\
\bottomrule
\end{tabularx}
\par\vspace{5pt}{\footnotesize\itshape পদ ও বাজেট blueprint/\allowbreak{}src/\allowbreak{}budget\_model.py (stage\_posts, totals ও scaleup ফাংশন) দ্বারা হিসাবকৃত; ধাপ ০ (phase0 ফাংশন) ৩৬ মাসের পরিকল্পনার আগে আসে ও আলাদাভাবে ব্যয়-নির্ধারিত; ধাপ ১ ও ২ হলো ৩৬ মাসের পরিকল্পনা; সম্প্রসারণ শর্তসাপেক্ষ ও অনুরোধবহির্ভূত। বছরগুলো নির্দেশক, বছর ১ হলো ২০২৭-২৮ অর্থবছর (মাস ১ জুলাই ২০২৭)।}
\par\vspace{3pt}{\scriptsize\color{gray} উৎস: BAPIL টেকনিক্যাল ব্লুপ্রিন্ট, সংস্করণ ১.০ (২০২৬)। CC BY 4.0।}
\end{minipage}
\end{document}
